% --- COLORES ---
\usepackage{xcolor}
% Definimos el azul institucional/cobalto
\definecolor{agebluelink}{HTML}{0F4C81}
% Un gris oscuro para subtítulos
\definecolor{darkgray}{HTML}{333333}

% --- ESTILO DE TÍTULOS (KOMA-SCRIPT) ---
% Aplicamos el color azul a los Capítulos (Actividades en tu documento)
\addtokomafont{chapter}{\color{agebluelink}\bfseries}
% Aplicamos el color azul a las Secciones (Tareas en tu documento)
\addtokomafont{section}{\color{agebluelink}\bfseries}
% Subsecciones en gris oscuro para crear contraste jerárquico
\addtokomafont{subsection}{\color{darkgray}\bfseries}

% Reducimos el enorme espacio en blanco que LaTeX suele dejar encima de los capítulos
\renewcommand*{\chapterheadstartvskip}{\vspace*{-\topskip}}
\renewcommand*{\chapterheadendvskip}{\vspace{1.5\baselineskip}}

% --- TABLAS ---
% Quarto ya usa booktabs por defecto, pero nos aseguramos de que el texto
% de las tablas tenga un poco más de aire (padding) para que respire.
\renewcommand{\arraystretch}{1.3}

% --- CALLOUTS (Cajas de aviso) ---
% Opcional: Si vas a usar ::: {.callout-note} en Quarto, el paquete tcolorbox
% que usa Quarto por detrás ya se encarga de que queden modernas. 
% Aquí no hace falta reescribir nada, Quarto lo hace genial por defecto en PDF.
